\begin{figure*}[!htbp] % Use figure* for spanning two columns
\begin{casestudypro}[title={Case Study. \textbf{(1)} Incorrect $\rightarrow$ \textbf{(2)} Mixed $\rightarrow$ \textbf{(3)} Correct}]
\textbf{Country:} Italy \\
\textbf{Rule-of-thumb:} It is expected to remove your shoes before entering someone's home. \\
\textbf{Story:} During a visit to her friend's house, Emily arrived 45 minutes late, carrying a bouquet of flowers for the host. Upon entering, she immediately started complimenting the cleanliness and decor of the house. Is what Emily did socially acceptable? \\
\textbf{Gold Label:} Neither \\ \\
\textbf{\textsc{Yi-1.5} (initial):} Yes. Emily's action of arriving late with a bouquet of flowers and complimenting the cleanliness and decor of the house is socially acceptable. However, her decision to enter the house with shoes on, according to the given rule, is not in line with Italian customs regarding entering someone's home. \raisebox{-0.1\height}{
\includegraphics[height=1em]{figure/icons/incorrect.png}} \\
\textbf{\textsc{SeaLLM-3} (initial):} Yes. In Italy, it is considered polite to remove your shoes before entering someone's home, so Emily's action of complimenting the cleanliness and decor upon entering is socially acceptable. \raisebox{-0.1\height}{
\includegraphics[height=1em]{figure/icons/incorrect.png}} \\
\textbf{\textsc{Yi-1.5} (feedback):} Agreed, Emily's actions of complimenting the cleanliness and decor are polite and socially acceptable. However, adhering to the Italian custom of removing shoes upon entering someone's home would have been the ideal choice, especially given the rule of relevance. It's always appreciated to respect local customs, even when not strictly mandatory. \\
\textbf{\textsc{SeaLLM-3} (feedback):} While Emily's compliments were polite, it would have been more socially acceptable for her to remove her shoes as per Italian customs. \\
\textbf{\textsc{Yi-1.5} (final):} No \raisebox{-0.1\height}{
\includegraphics[height=1em]{figure/icons/incorrect.png}} \\
\textbf{\textsc{SeaLLM-3} (final):} Neither \raisebox{-0.1\height}{\includegraphics[height=1em]{figure/icons/correct.png}} \\ \\
\textbf{Final decision:} Neither. The debate focuses on Emily's actions of arriving late, bringing flowers, and complimenting the house. \raisebox{-0.1\height}{\includegraphics[height=1em]{figure/icons/correct.png}}
\end{casestudypro}
\caption{Case study example of a debate with incorrect initial, mixed final and correct judge LLM decision.}

\end{figure*}